% TOP_PROBLEM: {"id": 3428, "title": "Atomicity of the poset of completary multifuncoids", "category_id": 7, "category": {"id": 7, "name": "topology", "display_name": "Topology", "description": "Properties preserved under continuous deformations.", "slug": "topology", "order_index": 7, "created_at": "2026-07-31T15:26:25.671Z"}, "status": "open"}
% FIRST_POSTED: null
% ENTRY_KIND: statement_only
% Imported from: batch04.tex; UnsolvedMath ID 3428
\documentclass[11pt]{article}
\usepackage[margin=25mm]{geometry}
\usepackage{fontspec}
\setmainfont{Latin Modern Roman}
\usepackage{amsmath,amssymb,mathtools}
\usepackage{unicode-math}
\setmathfont{Latin Modern Math}
\usepackage{xurl,hyperref}
\hypersetup{colorlinks=true,urlcolor=blue}
\setcounter{secnumdepth}{0}
\setlength{\parindent}{0pt}
\setlength{\parskip}{5pt}
\setlength{\emergencystretch}{3em}
\newcommand{\sref}[2]{\hyperlink{source:#1:#2}{[S#2]}}
\newcommand{\cataloguescope}[1]{\paragraph{Recorded target.}#1}
\begin{document}
% UnsolvedMath ID 3428; source catalogue code OPG-37386
\setcounter{section}{3427}
\section{Atomicity of the poset of completary multifuncoids}\label{3428}
{\small\sffamily UnsolvedMath ID 3428 · Topology}
\subsection{Definitions and mathematical statement}

\paragraph{Shared notation.}$\mathbb N=\{1,2,\ldots\}$, and $\mathbb Z,\mathbb Q,\mathbb R,\mathbb C$ denote the integers, rationals, reals and complex numbers. Any other conventions are specified locally.


For an index set $I$ and join-semilattices $L_i$ with bottom $0_i$, their product is ordered coordinatewise. Call $F\subseteq\prod_{i\in I}L_i$ completary when, for every $u,v$ in that product,
\[
 (u_i\vee v_i)_{i\in I}\in F
 \quad\Longleftrightarrow\quad
 \exists c:I\to\{0,1\}:\ (w_i)_{i\in I}\in F,
 \quad w_i=\begin{cases}u_i&c(i)=0,\\v_i&c(i)=1,\end{cases}
\]
and no tuple having a bottom coordinate belongs to $F$. This is the source's definition of a completary multifuncoid.

Order the specified collections of multifuncoids by set inclusion. The bottom is the empty multifuncoid. An atom is a nonempty element $A$ with no nonempty element strictly below it. Atomic means every nonempty element has an atom below it. Atomistic means every element is the least upper bound, in this poset, of all atoms below it; this specifies the join as a poset join rather than an unverified raw set union.
\paragraph{Statement recorded in the supplied source.}
For every set $X$ and index set $I$, is the inclusion poset of completary multifuncoids $F\subseteq\mathcal P(X)^I$ both atomic and atomistic? Here each coordinate lattice is $(\mathcal P(X),\subseteq,\cup,\varnothing)$. Both assertions and all finite or infinite index sets are included. \sref{3428}{2}
\subsection{Short English statement}
Do all nonempty completary multifuncoids on powersets contain an atom, and is each determined as the join of its atoms?
\subsection{Sources}
\begin{description}
\item[\hypertarget{source:3428:1}{S1}] ULAM.AI and the UnsolvedMath Contributors, UnsolvedMath: A Curated Collection of Open Mathematics Problems, record 3428 (OPG-37386). CC BY 4.0. \url{https://huggingface.co/datasets/ulamai/UnsolvedMath}.
\item[\hypertarget{source:3428:2}{S2}] Open Problem Garden, Atomicity of the poset of completary multifuncoids, OPG-37386. Original problem statement, full mathematical discussion, bibliography and comments. \url{https://www.openproblemgarden.org/op/atomicity_of_the_poset_of_multifuncoids}.
\end{description}
\label{end:3428}
\end{document}
